Die Aussage ist lokal, ohne Einschränkung sei daher

\mavergleichskette
{\vergleichskette
{ U
}
{ \subseteq }{  \R^n
}
{  }{
}
{    }{
}
{   }{
}
}
{}{}{}
versehen mit einer riemannschen Struktur durch Funktionen
\mathl{g_{ij}}{.} Die Voraussetzung besagt

\mavergleichskettedisp
{\vergleichskette
{ \left\langle \nabla_{\partial_i} \partial_j  ,  \partial_k   \right\rangle  + \left\langle  \partial_j  ,  \nabla_{\partial_i} \partial_k   \right\rangle
}
{ =} { \partial_i g_{jk}
}
{ } {
}
{ } {
}
{ } {}
}
{}{}{.}
Mit den Christoffelsymbolen formuliert ist

\mavergleichskettedisp
{\vergleichskette
{  \left\langle \nabla_{\partial_i} \partial_j   ,  \partial_k   \right\rangle
}
{ =} { \left\langle  \sum_{\ell = 1}^n \Gamma^\ell_{ij}  \partial_\ell   ,  \partial_k   \right\rangle
}
{ =} {   \sum_{\ell = 1}^n \Gamma^\ell_{ij}   g_{ \ell k}
}
{ } {
}
{ } {
}
}
{}{}{}
und

\mavergleichskettedisp
{\vergleichskette
{  \left\langle  \partial_j   ,  \nabla_{\partial_i} \partial_k   \right\rangle
}
{ =} { \left\langle  \partial_j   ,  \sum_{\ell = 1}^n \Gamma^\ell_{ik}  \partial_\ell     \right\rangle
}
{ =} {  \sum_{\ell = 1}^n \Gamma^\ell_{ik}   g_{ \ell j}
}
{ } {
}
{ } {
}
}
{}{}{,}
die Bedingung besagt also

\mavergleichskettedisp
{\vergleichskette
{   \sum_{\ell = 1}^n \Gamma^\ell_{ij}   g_{ \ell k} +\sum_{\ell = 1}^n \Gamma^\ell_{ik}   g_{ \ell j}
}
{ =} {   \partial_i g_{jk}
}
{ } {
}
{ } {
}
{ } {
}
}
{}{}{.}
Wir müssen

\mavergleichskettedisp
{\vergleichskette
{ D_V \left\langle  W  , Z  \right\rangle
}
{ =} { \left\langle \nabla_V W , Z  \right\rangle  + \left\langle  W  ,  \nabla_VZ  \right\rangle
}
{ } {
}
{ } {
}
{ } {
}
}
{}{}{}
für beliebige stetig differenzierbare Vektorfelder $V,W,Z$ zeigen, die wir jeweils als
\mathl{\sum_{i = 1}^nf_i \partial_i}{} mit stetig differenzierbaren Funktionen ansetzen können. Es sei zunächst

\mavergleichskette
{\vergleichskette
{ V
}
{ = }{ \partial_i
}
{  }{
}
{    }{
}
{   }{
}
}
{}{}{}
und

\mavergleichskettedisp
{\vergleichskette
{ W
}
{ =} { \sum_{j = 1}^nf_j \partial_j
}
{ } {
}
{ } {
}
{ } {
}
}
{}{}{}
und

\mavergleichskettedisp
{\vergleichskette
{ Z
}
{ =} { \sum_{k = 1}^n h_k \partial_k
}
{ } {
}
{ } {
}
{ } {
}
}
{}{}{.}
Das Skalarprodukt ist additiv in den Komponenten und $\nabla_{\partial_i}$ ist auch additiv,
da ein linearer Zusammenhang vorliegt. Ebenso ist

\mavergleichskette
{\vergleichskette
{ D_{\partial_i}
}
{ = }{ \partial_i
}
{  }{
}
{    }{
}
{   }{
}
}
{}{}{}
additiv. Deshalb können wir uns auf Felder der Form

\mavergleichskette
{\vergleichskette
{ W
}
{ = }{ f  \partial_j
}
{  }{
}
{    }{
}
{   }{
}
}
{}{}{}
und

\mavergleichskette
{\vergleichskette
{ Z
}
{ = }{ h  \partial_k
}
{  }{
}
{    }{
}
{   }{
}
}
{}{}{}
mit fixierten $j,k$ beschränken. Für die linke Seite ergibt sich

\mavergleichskettedisp
{\vergleichskette
{  \partial_i \left\langle f \partial_j  ,  h \partial_k   \right\rangle
}
{ =} { \partial_i \left(  fh g_{jk}  \right)
}
{ =} {  fg_{jk} \partial_i(h) + hg_{jk} \partial_i(f)   +  fh \partial_i \left(  g_{jk}   \right)
}
{ } {
}
{ } {
}
}
{}{}{.}
Nach
[[Differenzierbare Mannigfaltigkeit/R/Vektorbündel/Linearer Zusammenhang/Lokale Beschreibung/Fakt|Fakt]]
ist

\mavergleichskettedisp
{\vergleichskette
{ \nabla_{\partial_i} \left(  f \partial_j  \right)
}
{ =} {  \left(  \partial_i f  \right) \partial_j + f \nabla_{\partial_i }\partial_j
}
{ =} { \left(  \partial_i f  \right) \partial_j + f \left(  \sum_{ \ell = 1}^n \Gamma^\ell_{ij} \partial_\ell   \right)
}
{ } {
}
{ } {
}
}
{}{}{}
und

\mavergleichskettedisp
{\vergleichskette
{ \nabla_{\partial_i} \left(  h \partial_k  \right)
}
{ =} { \left(  \partial_i h  \right) \partial_k + h \nabla_{\partial_i }\partial_k
}
{ =} { \left(  \partial_i h  \right) \partial_k  + h \left(  \sum_{ \ell = 1}^n \Gamma^\ell_{ik} \partial_\ell   \right)
}
{ } {
}
{ } {
}
}
{}{}{.}
Daher ist

\mavergleichskettealign
{\vergleichskettealign
{  \left\langle  \nabla_{\partial_i} \left(  f \partial_j  \right)   ,  h \partial_k  \right\rangle
}
{ =} {  \left\langle  \left(  \partial_i f  \right) \partial_j  + f \left(  \sum_{ \ell = 1}^n \Gamma^\ell_{ij} \partial_\ell   \right)   ,  h \partial_k  \right\rangle
}
{ =} { h \left(   \partial_i f  \right) \left\langle   \partial_j  ,  \partial_k   \right\rangle +fh \sum_{ \ell = 1}^n   \Gamma_{ij}^\ell \left\langle          \partial_\ell   ,  \partial_k   \right\rangle
}
{ =} { h \left(   \partial_i f  \right)  g_{j k} +fh \sum_{ \ell = 1}^n   \Gamma_{ij}^\ell  g_{\ell k }
}
{ } {
}
}
{}
{}{}
und

\mavergleichskettealign
{\vergleichskettealign
{  \left\langle f \partial_j  ,  \nabla_{\partial_i} \left(  h \partial_k  \right)   \right\rangle
}
{ =} { \left\langle  f\partial_j  ,  \left(  \partial_i h \right) \partial_k  + h \left(  \sum_{ \ell = 1}^n \Gamma^\ell_{ik} \partial_\ell   \right)    \right\rangle
}
{ =} {   f \left(   \partial_i h  \right)  g_{j k} +fh \sum_{ \ell = 1}^n   \Gamma_{ik}^\ell  g_{\ell j }
}
{ } {
}
{ } {
}
}
{}
{}{.}
Die Summe der beiden Terme stimmt mit der linken Seite überein.

Im Richtungsfeld $V$ sind beide Seiten linear und auch mit der Multiplikation mit Funktionen verträglich.

 [[Kategorie:Latexseite]]
